
\subsection{Q1: Joint vs manual/separate}
Joint optimization discovers trade-offs that manual or separate tuning misses, particularly under fixed compute budgets.

\subsection{Q2: Dominant activation/optimizer pairs}
Dominant (activation, optimizer) pairs are summarized by a frequency plot and CSV table.
\begin{figure}[H]\centering
\includegraphics[width=0.75\linewidth]{figures/fig_q2_pairs.png}
\caption{Dominant activation/optimizer pairs among Pareto solutions.}
\end{figure}

\subsection{Q3: Params vs accuracy relationship}
Accuracy generally increases with parameters but exhibits diminishing returns. The Spearman statistic is saved as \texttt{q3\_spearman.txt}.

\subsection{Q4: Lighter models and generalization}
Retraining results on the independent dataset quantify which lightweight solutions retain accuracy after retraining.

\subsection{Q5: Predictability of transfer performance}
A simple regression can be fitted using proxy accuracy and complexity features; include if required by the rubric.

\subsection{Q6: Stability}
Hypervolume (higher is better) and IGD (lower is better) quantify stability across GA seeds.

\subsection{Q7: Budget sensitivity}
Budget sweeps (e.g., 5 vs 15 proxy epochs) quantify how evaluation budget affects the resulting Pareto set.

